\paragraph{坏模数的倍数密度与 window-\texorpdfstring{$6$}{6} 长根扇区的隐藏退化}
Dirichlet 扭曲坏模数沿整除链的传播，与 window-\(6\) 的 \(B_3\) 长根扇区在纤维多重度层上的残余对称，可进一步压缩为下列一组结构性结论。

\begin{corollary}[整除传播单独强迫坏模数集具有正下密度]\label{cor:xi-dirichlet-bad-moduli-positive-lower-density-by-divisibility}
\leanverified{paper\_xi\_dirichlet\_bad\_moduli\_positive\_lower\_density\_by\_divisibility}
沿用命题 \ref{prop:real-input-40-output-dirichlet-nonrh} 与结论 \ref{con:xi-time-part9l-dirichlet-twist-divisibility-monotonicity} 的记号，定义
\[
\mathcal B:=\{m\ge 2:\ \theta_m>\tfrac12\}.
\]
则 \(\mathcal B\neq\varnothing\)，并且若记
\[
m_0:=\min \mathcal B,
\]
则
\[
\underline d(\mathcal B)
:=
\liminf_{X\to\infty}\frac{\#(\mathcal B\cap [1,X])}{X}
\ge
\frac{1}{m_0}
>
0.
\]
换言之，即使不调用大奇素数模数的渐近展开，仅由“坏模数沿倍数链传播”这一整除刚性，也已足以推出坏模数集的正下密度。
\end{corollary}
\begin{proof}
由命题 \ref{prop:real-input-40-output-dirichlet-nonrh}，存在某个模数 \(m\ge 2\) 满足 \(\theta_m>\tfrac12\)，故 \(\mathcal B\neq\varnothing\)，从而最小元 \(m_0\) 存在。
再由结论 \ref{con:xi-time-part9l-dirichlet-twist-divisibility-monotonicity}，若 \(m_0\in\mathcal B\)，则其任意倍数 \(rm_0\)（\(r\ge 1\)）仍属于 \(\mathcal B\)。因此
\[
m_0\NN_{\ge 1}\subseteq \mathcal B.
\]
于是对任意 \(X\ge 1\)，有
\[
\#(\mathcal B\cap [1,X])\ge \left\lfloor \frac{X}{m_0}\right\rfloor.
\]
两边除以 \(X\) 并取 \(\liminf\)，即得
\[
\underline d(\mathcal B)\ge \frac{1}{m_0}.
\]
证毕。
\end{proof}

\begin{theorem}[window-\texorpdfstring{$6$}{6} dyadic 微态质量的 \texorpdfstring{$16+36+12$}{16+36+12} 精确三块分解]\label{thm:xi-window6-dyadic-microstate-exact-three-block-mass-split}
\leanverified{paper\_xi\_window6\_dyadic\_microstate\_exact\_three\_block\_mass\_split}
沿用
\[
\mathrm{Fold}^{\mathrm{bin}}_6:\{0,\dots,63\}\to X_6
\]
与 \(V_6\) 的记号。则
\[
2^6
=
\sum_{\substack{w\in X_6\\ w_6=1}} d_6^{\mathrm{bin}}(w)
+
\sum_{\substack{w\in X_6\\ w_6=0,\ V_6(w)\le 8}} d_6^{\mathrm{bin}}(w)
+
\sum_{\substack{w\in X_6\\ w_6=0,\ V_6(w)>8}} d_6^{\mathrm{bin}}(w)
=
16+36+12.
\]
等价地，全部 \(64\) 个 dyadic 微态在 \(\mathrm{Fold}^{\mathrm{bin}}_6\) 下被精确分解为三块总质量分别为 \(16\)、\(36\)、\(12\) 的互斥簇。进一步，在均匀微态基线下有
\[
\mathbb P(w_6=1)=\frac14,\qquad
\mathbb P(w_6=0,\ V_6(w)\le 8)=\frac{9}{16},\qquad
\mathbb P(w_6=0,\ V_6(w)>8)=\frac{3}{16}.
\]
\end{theorem}
\begin{proof}
由推论 \ref{cor:fold-bin-m6-fiber-hist}，
\[
\#\{w\in X_6:\ d_6^{\mathrm{bin}}(w)=2\}=8,\qquad
\#\{w\in X_6:\ d_6^{\mathrm{bin}}(w)=3\}=4,\qquad
\#\{w\in X_6:\ d_6^{\mathrm{bin}}(w)=4\}=9.
\]
另一方面，表 \ref{tab:fold6_bin_uplift_choice_collapse} 给出精确对应
\[
w_6=1 \Longleftrightarrow d_6^{\mathrm{bin}}(w)=2,
\]
\[
w_6=0,\ V_6(w)\le 8 \Longleftrightarrow d_6^{\mathrm{bin}}(w)=4,
\]
\[
w_6=0,\ V_6(w)>8 \Longleftrightarrow d_6^{\mathrm{bin}}(w)=3.
\]
因此三类纤维的总微态质量分别为
\[
8\cdot 2=16,\qquad 9\cdot 4=36,\qquad 4\cdot 3=12.
\]
又因为不同 \(w\) 的纤维两两不交且其并恰为 \(\{0,\dots,63\}\)，故
\[
64=16+36+12.
\]
最后两边除以 \(64\) 即得所示概率。证毕。
\end{proof}

\begin{theorem}[window-\texorpdfstring{$6$}{6} 的 \texorpdfstring{$A_2(+)$}{A2(+)} 与 \texorpdfstring{$A_2(-)$}{A2(-)} 长根扇区具有完全相同的纤维多重度谱]\label{thm:xi-window6-a2pm-identical-fiber-spectrum}
\leanverified{paper\_xi\_window6\_a2pm\_identical\_fiber\_spectrum}
沿用定理 \ref{thm:boundary-m10-b3c3-rigid-four-block-split} 的 \(B_3\) 根字典与分块记号。定义
\[
\Phi_{A_2(+)}
:=
\{\pm(e_1+e_2),\ \pm(e_1+e_3),\ \pm(e_2+e_3)\},
\]
\[
\Phi_{A_2(-)}
:=
\{\pm(e_1-e_2),\ \pm(e_1-e_3),\ \pm(e_2-e_3)\},
\]
\[
X_6^{A_2(+)}:=\iota_B^{-1}\bigl(\Phi_{A_2(+)}\bigr),
\qquad
X_6^{A_2(-)}:=\iota_B^{-1}\bigl(\Phi_{A_2(-)}\bigr).
\]
则两块的纤维多重度多重集完全一致，且都等于
\[
\{4,4,3,3,2,2\}.
\]
等价地，若定义生成多项式
\[
P_{A_2(\pm)}(t):=\sum_{w\in X_6^{A_2(\pm)}} t^{d_6^{\mathrm{bin}}(w)},
\]
则
\[
P_{A_2(+)}(t)=P_{A_2(-)}(t)=2t^4+2t^3+2t^2.
\]
\end{theorem}
\begin{proof}
由表 \ref{tab:boundary_uplift_m10_dictionary} 可直接读出
\[
X_6^{A_2(+)}
=
\{\texttt{000000},\texttt{010100},\texttt{100010},
\texttt{101010},\texttt{010001},\texttt{010101}\},
\]
\[
X_6^{A_2(-)}
=
\{\texttt{101000},\texttt{100100},\texttt{010010},
\texttt{001010},\texttt{001001},\texttt{000101}\}.
\]
另一方面，推论 \ref{cor:fold-bin-m6-fiber-hist} 与表 \ref{tab:fold6_bin_uplift_choice_collapse} 给出三档精确分类：
\[
w_6=0,\ V_6(w)\le 8 \Longleftrightarrow d_6^{\mathrm{bin}}(w)=4,
\]
\[
w_6=0,\ V_6(w)>8 \Longleftrightarrow d_6^{\mathrm{bin}}(w)=3,
\]
\[
w_6=1 \Longleftrightarrow d_6^{\mathrm{bin}}(w)=2.
\]
其中在 \(X_6^{A_2(+)}\) 中，
\[
\texttt{000000},\texttt{010100}\in\{w_6=0,\ V_6(w)\le 8\},
\]
\[
\texttt{100010},\texttt{101010}\in\{w_6=0,\ V_6(w)>8\},
\]
\[
\texttt{010001},\texttt{010101}\in\{w_6=1\}.
\]
故其多重度多重集为
\[
\{4,4,3,3,2,2\}.
\]
同理，在 \(X_6^{A_2(-)}\) 中，
\[
\texttt{101000},\texttt{100100}\in\{w_6=0,\ V_6(w)\le 8\},
\]
\[
\texttt{010010},\texttt{001010}\in\{w_6=0,\ V_6(w)>8\},
\]
\[
\texttt{001001},\texttt{000101}\in\{w_6=1\},
\]
故其多重度多重集同样为
\[
\{4,4,3,3,2,2\}.
\]
于是两块的生成多项式均为
\[
2t^4+2t^3+2t^2.
\]
证毕。
\end{proof}

\begin{corollary}[任何仅依赖纤维多重度的统计量都无法区分 \texorpdfstring{$A_2(+)$}{A2(+)} 与 \texorpdfstring{$A_2(-)$}{A2(-)}]\label{cor:xi-window6-a2pm-multiplicity-blind-statistics}
\leanverified{paper\_xi\_window6\_a2pm\_multiplicity\_blind\_statistics}
沿用定理 \ref{thm:xi-window6-a2pm-identical-fiber-spectrum} 的记号。对任意函数
\[
\Phi:\{2,3,4\}\to \RR,
\]
都有
\[
\sum_{w\in X_6^{A_2(+)}}\Phi\!\bigl(d_6^{\mathrm{bin}}(w)\bigr)
=
\sum_{w\in X_6^{A_2(-)}}\Phi\!\bigl(d_6^{\mathrm{bin}}(w)\bigr).
\]
特别地，对任意 \(q\in\RR\)，有
\[
\sum_{w\in X_6^{A_2(+)}}\bigl(d_6^{\mathrm{bin}}(w)\bigr)^q
=
\sum_{w\in X_6^{A_2(-)}}\bigl(d_6^{\mathrm{bin}}(w)\bigr)^q
=
2(2^q+3^q+4^q).
\]
\end{corollary}
\begin{proof}
由定理 \ref{thm:xi-window6-a2pm-identical-fiber-spectrum}，两块的纤维多重度多重集都等于
\[
\{4,4,3,3,2,2\}.
\]
因此对任意只依赖多重度数值的标量函数 \(\Phi\)，两侧求和都仅由同一个多重集决定，故必相等。取 \(\Phi(d)=d^q\) 即得第二式。证毕。
\end{proof}

\begin{theorem}[纤维大小场在 \texorpdfstring{$W(B_3)$}{W(B3)}-不变子空间中的最优两轨道压缩]\label{thm:xi-window6-weyl-invariant-best-two-orbit-compression}
沿用推论 \ref{cor:fold6-weyl-two-orbit-compression} 的分块记号，记
\[
d_6^{\mathrm{cyc}}:=d_6^{\mathrm{bin}}\big|_{X_6^{\mathrm{cyc}}}:X_6^{\mathrm{cyc}}\to\{2,3,4\}.
\]
在实向量空间 \(\RR^{X_6^{\mathrm{cyc}}}\) 上取均匀内积
\[
\langle f,g\rangle
:=
\frac{1}{18}\sum_{w\in X_6^{\mathrm{cyc}}}f(w)g(w),
\]
并定义 \(W(B_3)\)-不变子空间
\[
\mathcal W
:=
\{f:X_6^{\mathrm{cyc}}\to\RR:\ f\ \text{在}\ X_6^{\mathrm{short}}\ \text{与}\ X_6^{\mathrm{long}}\ \text{上分别常值}\}.
\]
则成立：
\begin{enumerate}
\item \(d_6^{\mathrm{cyc}}\) 在 \(\mathcal W\) 上的唯一正交投影 \(\Pi_{\mathcal W}d_6^{\mathrm{cyc}}\) 为
\[
(\Pi_{\mathcal W}d_6^{\mathrm{cyc}})(w)
=
\begin{cases}
\frac{11}{3},& w\in X_6^{\mathrm{short}},\\[4pt]
3,& w\in X_6^{\mathrm{long}}.
\end{cases}
\]
\item 最小均方误差恰为
\[
\inf_{f\in\mathcal W}
\frac{1}{18}\sum_{w\in X_6^{\mathrm{cyc}}}\bigl(d_6^{\mathrm{cyc}}(w)-f(w)\bigr)^2
=
\frac{17}{27}.
\]
\item 一致范数下的最佳 \(W(B_3)\)-不变逼近误差恰为
\[
\inf_{f\in\mathcal W}\|d_6^{\mathrm{cyc}}-f\|_{\infty}=1,
\]
且唯一极小元为常函数
\[
f\equiv 3.
\]
\end{enumerate}
\leanverified{paper\_xi\_window6\_weyl\_invariant\_best\_two\_orbit\_compression}
\end{theorem}
\begin{proof}
由表 \ref{tab:fold6_b3c3_root_dictionary_b3} 与表 \ref{tab:fold_bin_gauge_m6_fibers}，六个短根原像的多重度多重集为
\[
\{4,4,4,4,4,2\}.
\]
另一方面，由定理 \ref{thm:xi-window6-a2pm-identical-fiber-spectrum}，
\[
X_6^{\mathrm{long}}=X_6^{A_2(+)}\sqcup X_6^{A_2(-)}
\]
且两块各自的多重度多重集都为 \(\{4,4,3,3,2,2\}\)。故十二个长根原像的多重度多重集为
\[
\{4,4,4,4,3,3,3,3,2,2,2,2\}.
\]

由推论 \ref{cor:fold6-weyl-two-orbit-compression}，\(\mathcal W\) 正是“短根块常值、长根块常值”的二维子空间。因此其正交投影由两条轨道上的平均值给出。短根平均值为
\[
\frac{5\cdot 4+2}{6}=\frac{11}{3},
\]
长根平均值为
\[
\frac{4\cdot 4+4\cdot 3+4\cdot 2}{12}=3.
\]
这就得到第一条。

对第二条，短根块的最小平方误差贡献为
\[
\frac{1}{18}\left(
5\left(4-\frac{11}{3}\right)^2+\left(2-\frac{11}{3}\right)^2
\right)
=
\frac{1}{18}\left(5\cdot \frac{1}{9}+\frac{25}{9}\right)
=
\frac{5}{27},
\]
长根块的最小平方误差贡献为
\[
\frac{1}{18}\left(
4(4-3)^2+4(3-3)^2+4(2-3)^2
\right)
=
\frac{8}{18}
=
\frac{4}{9}.
\]
两者相加即得
\[
\frac{5}{27}+\frac{4}{9}=\frac{17}{27}.
\]

最后考察一致范数。任取 \(f\in\mathcal W\)，则存在 \(a,b\in\RR\) 使
\[
f|_{X_6^{\mathrm{short}}}\equiv a,
\qquad
f|_{X_6^{\mathrm{long}}}\equiv b.
\]
于是
\[
\sup_{w\in X_6^{\mathrm{short}}}|d_6^{\mathrm{cyc}}(w)-f(w)|
=
\max\{|2-a|,|4-a|\},
\]
其最小值为 \(1\)，且唯一在 \(a=3\) 时取得；
\[
\sup_{w\in X_6^{\mathrm{long}}}|d_6^{\mathrm{cyc}}(w)-f(w)|
=
\max\{|2-b|,|3-b|,|4-b|\},
\]
其最小值同样为 \(1\)，且唯一在 \(b=3\) 时取得。故
\[
\|d_6^{\mathrm{cyc}}-f\|_\infty
=
\max\left\{
\max\{|2-a|,|4-a|\},
\max\{|2-b|,|3-b|,|4-b|\}
\right\}
\ge 1,
\]
并且等号当且仅当 \(a=b=3\) 时成立。于是最优一致逼近误差为 \(1\)，唯一极小元为常函数 \(f\equiv 3\)。证毕。
\end{proof}

\noindent
这些结论说明：Dirichlet 扭曲侧的平方根门槛失效一旦出现，便会沿倍数链产生具有正下密度的坏模数族；而在 window-\(6\) 的 \(B_3\) 离散根字典中，虽然纤维大小场严格打破了 Weyl 轨道均匀性，长根层内部却仍保留一条精确的 \(A_2(+)/A_2(-)\) 多重度退化。于是“Weyl 破缺”与“残余二块等谱”并非矛盾，而是在不同粗细尺度上同时成立的两级结构。

\endinput
